\section{附录：常见实践问题}

在训练深度 RL 时无可避免地会遇到一些反复出现的问题，这里记录几个讲座中提到或合理推断的 case：

\begin{enumerate}
    \item \textbf{为何策略出现 sudden collapse}: 这通常发生在 importance sampling weight 太高或 clipping 超出 1.2 之后。建议每训练 epoch 后计算 `std(ratio)` 和 `max(ratio)`，若均值超过 0.12 则退回上个 checkpoint。
    \item \textbf{为何 replay buffer 中存在重复轨迹}: 原因通常是高频重置（reset 从相近状态开始）。解决方法是在 buffer 中插入 `prioritized eviction`，优先清除 `similarity > 0.9` 的 transition 对。
    \item \textbf{如何控制 entropy decay}: 在 SAC 中 entropy 系数 $\alpha$ 不应固定，而是让 value head 估计 entropy target（例如 0.2），并用 `log_alpha` 自动调节。
    \item \textbf{如何对抗 reward hacking}: 设立 `constraint auxiliary rewards`，例如 `-0.2` 罚金 for hitting unnatural states，并在 reward shaping 中使用 potential-based theorem 保证 optimal policy 不变。
    \item \textbf{如何保持 training trace}: 使用 structured logging 记录 `config, seed, buffer stats, loss curves`，并将这些记录上传至 centralized experiment dashboard，便于 future debugging。
    \item \textbf{何时应切换算法}: 如果 training variance 超出 checkpoint window（例如 10 连续 episodes variance > 50\% of mean），可以将 job 转成 PPO + SAC dual training，让 PPO 作为 anchor policy。
\end{enumerate}

\begin{warningbox}{别把附录当作 checklist}
这些问题只是常见 heuristics，不能替代具体环境的 ablation。即便在相似的 env 中，也应该验证 `variance`、`kl` 与 `reward` 的纵向变化，并据此微调 thresholds。
\end{warningbox}

\end{document}
